\SourceLeaf{405}%
wobei \(c\) den Unterschied zwischen den Radien des ersten und des zweiten Kreises und \(c'\) den Unterschied zwischen den Radien des ersten und des dritten Kreises bezeichnet.

Wir haben die folgenden Gleichungen
\[
\begin{aligned}
xx+yy&=zz\\
(x-a)^2+(y-b)^2&=(z-c)^2\\
(x-a')^2+(y-b')^2&=(z-c')^2
\end{aligned}
\]
Setzt man
\[
\begin{aligned}x&=z\cos\varphi\\y&=z\sin\varphi\end{aligned}
\]
so erhält man aus den vorigen Gleichungen
\[
\begin{aligned}
aa+bb-cc&=2az\cos\varphi+2bz\sin\varphi-2cz\\
a'a'+b'b'-c'c'&=2a'z\cos\varphi+2b'z\sin\varphi-2c'z
\end{aligned}
\]
Dividiert man die erste Gleichung durch \(a\cos\varphi+b\sin\varphi-c\), die zweite durch \(a'\cos\varphi+b'\sin\varphi-c'\) und zieht sie dann voneinander ab, erhält man
\[
\frac{a'a'+b'b'-c'c'}{a'\cos\varphi+b'\sin\varphi-c'}-\frac{aa+bb-cc}{a\cos\varphi+b\sin\varphi-c}=0
\]
oder, wenn wir schreiben
\[
\begin{aligned}a'a'+b'b'-c'c'&=A'\\aa+bb-cc&=A\end{aligned}
\]
ergibt sich
\[
(A'a-Aa')\cos\varphi+(A'b-Ab')\sin\varphi=A'c-Ac'
\]
Setzen wir nun
\[
\begin{aligned}A'a-Aa'&=R\cos M\\A'b-Ab'&=R\sin M\\A'c-Ac'&=N\end{aligned}
\]
so geht unsere Gleichung über in
\[
R\cos(\varphi-M)=N
\]
Darin ist außer \(\varphi\) alles gegeben. Man erhält daraus zwei Werte für \(\varphi\), zwischen denen man je nach der gewünschten Art der Berührung wählen muss.
